\section{实验结果}
\label{sec:results}

本文方法的整体流程如 \Cref{alg:pd-cms} 所示，实现细节见 \Cref{sec:implementation}。我们在三类平台上进行测试：用于扩展性基准的桌面工作站、用于单机对比的服务器，以及用于超大规模实验的超算集群。

\begin{algorithm}[H]
  \caption{多重划分 CMS}
  \label{alg:pd-cms}
  \begin{algorithmic}[1]
    \REQUIRE $nev_{\text{blk}},nev_{\text{target}}$
    \STATE \COMMENT{\textbf{并行运行}以获取 $\mathbf{S}_{\mathcal{I}}^p,\mathbf{S}_{\mathcal{I}}^d$}
    \FOR{each substructure in $\mathcal{P}^p$ and $\mathcal{P}^d$}
      \STATE 求解子结构特征值问题 \cref{eq:sub-eig}
    \ENDFOR
    \STATE \COMMENT{\textbf{并行运行}以获取 $\mathbf{S}_b^p,\mathbf{S}_b^d$}
    \FOR{each substructure in $\mathcal{P}^p$ and $\mathcal{P}^d$}
      \STATE 求解子结构边界问题 \cref{eq:ours-response-solution} 并组装
    \ENDFOR
    \STATE $\mathbf{S} \leftarrow [\mathbf{S}^p_\mathcal{I},\mathbf{S}^p_b,\mathbf{S}_{\mathcal{I}}^d,\mathbf{S}_{b}^d]$
    \STATE 使用 \cref{alg:adaptive-shifting} 的自适应位移过程求解最终缩减系统
  \end{algorithmic}
\end{algorithm}

我们首先在标准二维矩形区域（四边形网格）上验证方法的基础数值行为。随后在两类问题上进行对比：三角网格上的拉普拉斯特征值问题，以及四面体网格上的弹性特征值问题。进一步地，我们在超算集群上验证了超大规模模型的效率，并展示了局部网格改动下的快速局部更新能力（CMS 的关键优势之一）。

\begin{table*}[H]
  \caption{\textbf{实验所用模型统计。} $q_t$ 表示需计算的全局特征模态数量；* 表示变量。}
  \label{tab:statistics}
  \centering
  \begin{tabular}{|c|c|c|c|c|}
    \hline
    平台 & 模型 & \#V & \#E & $q_t$  \\
    \hline
    \multirow{5}*{\makecell{桌面工作站\\AMD Ryzen 9 7900X 12-Core Processor @ 3.6GHz\\128 GiB RAM}}
    & rectangle（\Cref{fig:scaling}） & * & * & * \\
    & PCB（\Cref{fig:pcb-local-update}） & 320K & 642K & 100 \\
    & PCB modified（\Cref{fig:pcb-local-update}） & 319K & 638K & 100 \\
    & microstructure（\Cref{fig:micro-local-update}） & 321K & 105K & 100 \\
    & microstructure modified（\Cref{fig:micro-local-update}） & 333K & 111K & 100 \\
    \hline
    \multirow{4}*{\makecell{计算服务器\\2 $\times$ Intel Xeon Gold 6254 @ 3.10GHz\\1024 GiB RAM}}
    & Frame（\Cref{fig:comp-timing-vs-nev}） & 137K & 274K & * \\
    & Breaker（\Cref{fig:laplacian-aa-error-time}） & 3.3M & 6.6M & 500 \\
    & Cat（\Cref{fig:laplacian-metis-error-time}） & 765K & 1.5M & 500 \\
    & Wing（\Cref{fig:elasticity-aa-error-time}） & 40K & 140K & 200 \\
    \hline
    \makecell{超算集群（19 节点）\\AMD 2$\times$7H12 64-core @ 2.6GHz, 512GB RAM\\100Gb IB 高速网络}
    & Falcon 9（\Cref{fig:teaser}） & 10M & 30M & 3000 \\
    \hline
  \end{tabular}
\end{table*}

\subsection{精度与性能评估}
\label{sec:evaluation}

\paragraph{误差度量}
我们根据是否具备高精度基准解，采用不同精度指标评估特征对 $(\Lambda,\Phi)$。当问题规模允许计算参考特征值 $\lambda_q^*$（由直接法获得，如 Spectra~\cite{Qiu2015Spectra} + CHOLMOD~\cite{2008ChenCholmod}）时，使用相对特征值误差：
\begin{equation}
  \epsilon_{\text{ev}} = \sqrt{\sum_{q=1}^{N_e} \left( \frac{|\lambda_q - \lambda_q^*|}{\lambda_q^*} \right)^2},
  \label{eq:error-ev}
\end{equation}
其中 $N_e$ 为已计算特征值数量。

\paragraph{性能评估：误差--时间曲线}
我们将误差（$\epsilon_{\text{ev}}$）对总时间（$t_{\text{cms}}$）绘制成曲线。曲线越靠左下，表示在同等时间预算下误差更低。
\subsection{可扩展性}
\label{sec:scaling}
可扩展性是 HPC 的核心指标。我们使用\textbf{强扩展}（固定问题规模，增加处理器数量）与\textbf{弱扩展}（问题规模与处理器数量同比增长）两类标准指标。现代集群通常采用 MPI（跨节点）+ OpenMP（节点内）两级并行模型，总并行度为 $p_c\times p_f$。

CMS 与该架构天然匹配：子结构分解提供粗粒度、易并行的任务划分；每个子结构内的局部特征值求解主要依赖局部数据，可有效利用 OpenMP 并行。

\subsubsection{MPI 环境模拟与性能指标}
受限于实验资源，我们采用与 SMPI 类似的\textbf{并行环境模拟}策略~\cite{Degomme2017SMPI}：顺序执行 $m$ 个子结构任务并记录每个任务的时间 $T_i$ 与峰值内存 $M_i$，据此估计并行指标：
\begin{itemize}
  \item \textbf{估计并行总时间}：$T_{\text{parallel}}\approx\max_{i=1,\dots,m}(T_i)$；
  \item \textbf{峰值节点内存}：$M_{\text{peak}}=\max_{i=1,\dots,m}(M_i)$；
  \item \textbf{负载均衡（时间/内存）}：采用平均值与最大值之比，越接近 1 表示均衡越好。
\end{itemize}

\subsubsection{实验设置}
强扩展实验固定网格规模；弱扩展实验按节点数同比放大网格。为公平比较，统一使用阈值 $\epsilon_{\text{ev}}<\tau=10^{-4}$，并通过二分策略搜索恰好满足阈值的关键内部模态数 $q_{\mathcal{I}}^\star$。

本文方法采用两套划分：$m^p$ 与 $m^d=m^p+1$，总节点数为 $2m^p+1$。传统 CB-CMS 使用相同总节点数并直接划分为 $2m^p+1$ 个子结构。

\subsubsection{结果}
$q_{\mathcal{I}}^\star$ 见 \Cref{fig:critical-qb}。\textcolor{ConclusionEvidence}{受益于相位补足，本文方法仅需较少子结构模态，而传统 CMS 往往需要数百个。}

\Cref{fig:scaling} 显示本文方法的时间扩展趋势与 CB-CMS 接近，但总体速度快约两个数量级；峰值内存扩展性显著更好。该结论与我们的分析一致：传统 CMS 的峰值内存受 Schur 补稠密界面矩阵支配并随界面规模近似二次增长；MP-CMS 消除了该瓶颈，内存占用主要受稀疏子结构矩阵控制。

\textcolor{IsConclusion}{此外，分解结果清楚表明，传统 CMS 的串行界面步骤（Schur 补）会随问题规模增大而在内存占比上持续上升，这正是其内存可扩展性变差的直接原因。}

\Cref{fig:scaling-balance} 显示，本文方法在强、弱扩展场景下的时间与内存均衡性均约为 $0.9$。传统 CB-CMS 方法则明显下降，特别是内存均衡性接近 $0.6$。
\begin{figure}[H]
  \centering
  \includegraphics[width=\linewidth]{critical-qb/combined_scaling.pdf}
  \caption{\textbf{子结构特征模态关键数量。} 本文方法所需 $q_{\mathcal{I}}^\star$ 明显小于传统 CMS。}
  \label{fig:critical-qb}
\end{figure}

\begin{figure}[H]
  \centering
  \includegraphics[width=\linewidth]{scaling_aggregation_new/combined_plots/combined_strong_weak_total_breakdown.pdf}
  \includegraphics[width=\linewidth]{scaling_aggregation_new/combined_plots/combined_strong_weak_peak_memory_line.pdf}
  \caption{\textbf{强扩展与弱扩展下的时间/内存代价。} 本文方法时间成本约低两个数量级，且内存成本和内存扩展性均显著优于传统 CMS。}
  \label{fig:scaling}
\end{figure}

\begin{figure}[H]
  \centering
  \includegraphics[width=\linewidth]{scaling_aggregation_new/combined_plots/combined_strong_weak_balance_line.pdf}
  \caption{\textbf{并行均衡性（强扩展与弱扩展）。} 本文方法在时间与内存维度均表现出更优均衡性。}
  \label{fig:scaling-balance}
\end{figure}

\subsection{方法对比}
本节将本文方法（\textbf{MP-CMS}）与三种基线进行对比：\textbf{CB-CMS}（单划分+完整 Schur 补）、\textbf{CMS-IMR}（单划分+界面模态缩减）和 \textbf{AMLS}~\cite{Gao2008AMLS}。

\subsubsection{误差--时间分析}
\paragraph{轴对齐划分下的拉普拉斯问题}
\Cref{fig:laplacian-aa-error-time} 显示：CB-CMS 因 Schur 补稠密计算具有较高初始开销；CMS-IMR 与 AMLS 虽降低了部分代价，但仍需较重的基构建时间。\textcolor{ConclusionEvidence}{MP-CMS 在相同预算下达到最低误差，较基线普遍优 $1$--$2$ 个数量级。}

\paragraph{METIS 划分下的拉普拉斯问题}
在不规则子域场景中（\Cref{fig:laplacian-metis-error-time}），尽管 METIS 对 CMS-IMR 与 AMLS 亦有帮助，MP-CMS 仍保持优势。并且，若不加入互补子结构（PD-NC），精度显著退化，验证了交点区域相位补足机制的必要性。

\paragraph{轴对齐划分下的弹性问题}
\Cref{fig:elasticity-aa-error-time} 表明，三维弹性问题中 MP-CMS 同样优于 CB-CMS、CMS-IMR 与 AMLS。\textcolor{ConclusionEvidence}{在相同时间下可提升约 3 个数量级精度，或在同精度下节省约 3 个数量级时间。}

\begin{figure}[H]
    \centering
    \begin{subfigure}{0.8\linewidth}
        \includegraphics[width=0.189\linewidth]{comparison-figures/breaker/primal-partition.png}
        \includegraphics[width=0.189\linewidth]{comparison-figures/breaker/dual-partition.png}
        \vrule\
        \includegraphics[width=0.189\linewidth]{comparison-figures/breaker/eig-1.png}
        \includegraphics[width=0.189\linewidth]{comparison-figures/breaker/eig-25.png}
        \includegraphics[width=0.189\linewidth]{comparison-figures/breaker/eig-47.png}
        \subcaption{左：primal 与 dual 划分；右：第 1、25、47 阶全局模态。}
    \end{subfigure}
    \vspace{3mm}
    \begin{subfigure}{\linewidth}
        \includegraphics[width=\linewidth]{comparison-figures/breaker/breaker_comparison.pdf}
    \end{subfigure}
    \caption{\textbf{轴对齐划分下拉普拉斯问题对比。} 本文方法在该场景下优于其他 CMS 类方法。}
    \label{fig:laplacian-aa-error-time}
\end{figure}

\begin{figure}[H]
    \centering
    \begin{subfigure}{\linewidth}
      \includegraphics[width=0.189\linewidth]{comparison-figures/cat/primal-partition.png}
      \includegraphics[width=0.189\linewidth]{comparison-figures/cat/dual-partition.png}
      \vrule\
      \includegraphics[width=0.189\linewidth]{comparison-figures/cat/eig-1.png}
      \includegraphics[width=0.189\linewidth]{comparison-figures/cat/eig-25.png}
      \includegraphics[width=0.189\linewidth]{comparison-figures/cat/eig-47.png}
      \subcaption{左：不同边权下生成的 METIS primal/dual 划分；右：第 1、25、47 阶全局模态。}
      \label{fig:metis-pd-partition}
    \end{subfigure}
    \vspace{3mm}
    \begin{subfigure}{\linewidth}
      \centering
      \includegraphics[width=0.333\linewidth]{comparison-figures/cat/comp-2.png}%
      \includegraphics[width=0.206\linewidth]{comparison-figures/cat/comp-3.png}%
      \includegraphics[width=0.306\linewidth]{comparison-figures/cat/comp-4.png}%
      \includegraphics[width=0.154\linewidth]{comparison-figures/cat/comp-5.png}
      \subcaption{当 $\mathcal{P}^p$ 与 $\mathcal{P}^d$ 界面相交时，需要引入互补子结构（蓝色）以恢复相位完备性。}
      \label{fig:metis-c-partition}
    \end{subfigure}
    \vspace{3mm}
    \begin{subfigure}{\linewidth}
      \includegraphics[width=\linewidth]{comparison-figures/cat/cat_comparison.pdf}
    \end{subfigure}
    \caption{\textbf{METIS 划分下拉普拉斯问题对比。} 本文方法在相同时间预算下达到更低误差；无互补子结构时（PD-NC）误差明显升高。}
    \label{fig:laplacian-metis-error-time}
\end{figure}

\begin{figure}[H]
  \centering
  \begin{subfigure}{\linewidth}
    \centering
    \includegraphics[height=0.52\linewidth]{comparison-figures/board/primal-partition.png}\hfill
    \includegraphics[height=0.52\linewidth]{comparison-figures/board/dual-partition.png}\hfill
    \vrule\
    \includegraphics[height=0.52\linewidth]{comparison-figures/board/eig-7.png}\hfill
    \includegraphics[height=0.52\linewidth]{comparison-figures/board/eig-103.png}\hfill
    \includegraphics[height=0.52\linewidth]{comparison-figures/board/eig-199.png}
    \subcaption{左：4、5 子结构轴对齐划分；右：第 7、103、199 阶全局模态。}
  \end{subfigure}
  \vspace{3mm}
  \begin{subfigure}{\linewidth}
    \includegraphics[width=\linewidth]{comparison-figures/board/board_comparison.pdf}
  \end{subfigure}
  \caption{\textbf{轴对齐划分下弹性问题对比。} 本文方法在精度与效率上均显著优于对比方法。}
  \label{fig:elasticity-aa-error-time}
\end{figure}

\subsubsection{时间 vs. 准确模态数量}
为进一步量化实用价值，我们统计在给定阈值 $\epsilon_{\text{target}}$ 下达到目标准确模态数 $q_t$ 所需总时间 $t_{\text{cms}}$。
\textcolor{ConclusionEvidence}{如 \Cref{fig:comp-timing-vs-nev} 所示，MP-CMS 在绝大多数目标 $q_t$ 上均快于 CB-CMS、CMS-IMR 与 AMLS，且在大规模模态需求时优势更加显著。}
\begin{figure}[H]
 \centering
 \includegraphics[width=\linewidth]{time-nev/nev_vs_time_plot_sparse.pdf}
 \caption{\textbf{整体效率对比：时间 vs. 准确模态数量。} 每个点对应一个 $q_{\mathcal{I}}$ 设定下的独立运行。}
 \label{fig:comp-timing-vs-nev}
\end{figure}

\subsection{可扩展能力}
\subsubsection{可重用性}
与全局整体式求解器在局部改动后需整体重算不同，本文方法利用子结构模块化实现高效局部更新。当子结构 $\Omega_i$ 被修改时，仅需更新：（1）$\Omega_i$ 的局部特征模态；（2）与 $\Omega_i$ 相关的 SE-free 界面模态。

\Cref{fig:pcb-local-update} 展示了 PCB 模型拓扑删减案例，局部更新较完全重算实现约 $2\times$ 加速。\Cref{fig:micro-local-update} 展示了微结构弹性案例，参数化修改后，局部更新较完整 MP-CMS 重算加速 $5.4\times$，较 Spectra+CHOLMOD 仍有 $2.24\times$ 提升。

\begin{figure}[H]
  \centering
  \begin{tabular}{@{} m{1.5em} c c @{} }
    & \textbf{Primal} & \textbf{Dual} \\[1ex]

    \rotatebox{90}{\textbf{Origin}} &
    \includegraphics[width=0.45\linewidth, valign=m]{local_update/pcb_origin_primal.png} &
    \includegraphics[width=0.45\linewidth, valign=m]{local_update/pcb_origin_dual.png} \\

    \rotatebox{90}{\textbf{Modified}} &
    \includegraphics[width=0.45\linewidth, valign=m]{local_update/pcb_updated_primal.png} &
    \includegraphics[width=0.45\linewidth, valign=m]{local_update/pcb_updated_dual.png} \\
  \end{tabular}

  \includegraphics[width=\linewidth]{local_update/PCB_np4_neig100/benchmark_comparison.pdf}
  \caption{\textbf{PCB 模型：划分与局部更新性能。} 局部更新将总时间由 20.02\,s 降至 15.60\,s。}
  \label{fig:pcb-local-update}
\end{figure}

\begin{figure}[H]
  \centering
  \begin{tikzpicture}
    \node[anchor=north west, inner sep=0] (img) at (0,0) {%
      \includegraphics[width=\linewidth]{local_update/micro_long/partition-and-res.png}%
    };

    \pgfmathsetlengthmacro{\imgheight}{0.374*\linewidth}
    \node[anchor=west, font=\footnotesize\bfseries] at ([xshift=0.30\linewidth, yshift=-0.33*\imgheight]img.north west) {原始网格的划分};
    \node[anchor=west, font=\footnotesize\bfseries] at ([xshift=0.30\linewidth, yshift=-0.7*\imgheight]img.north west) {修改后网格的划分};
    \node[anchor=north west, font=\footnotesize\bfseries] at ([xshift=0.30\linewidth, yshift=-0.3mm]img.south west) {特征向量对比};
    \draw[red, line width=.5pt] (img.north west) rectangle ([xshift=0.05\linewidth]img.south west);
  \end{tikzpicture}
  \includegraphics[width=\linewidth]{local_update/micro_long/benchmark_neig100_comparison.pdf}
  \caption{\textbf{微结构模型：划分与局部更新性能。} 局部更新相较完整重算约加速 $2.24\times$。}
  \label{fig:micro-local-update}
\end{figure}

\subsubsection{超大规模模型性能}
为验证分布式环境下的实际可行性，我们在 19 节点 MPI/OpenMP 集群上测试 Falcon 9 壳模型（1000 万顶点，3000 万 DoFs），primal/dual 分别划分为 9/10 个子结构，每个子结构取 100 个局部模态，目标求解前 3000 阶全局模态。

\begin{figure}[H]
  \centering
  \includegraphics[width=\linewidth]{teaser/teaser.pdf}
  \caption{\textbf{Falcon 9 模型上的大规模 MPI 并行计算。} 本文方法扩展为基于交错多重划分的 MP-CMS。总耗时 88 分钟，其中通信仅 4.3 秒；在该规模下，传统 CB-CMS 会因界面稠密矩阵导致内存需求约 2.82 TB/节点而不可行。对比地，SLEPc+MUMPS 在更宽松阈值 $3\times10^{-2}$ 下约 120 分钟仅收敛 949 阶；本文方法则完成 3000 阶并保持多数频段误差在 $10^{-6}$ 量级。}
  \label{fig:teaser}
\end{figure}

\begin{table}[H]
  \centering
  \caption{\textbf{Falcon 9 超算实验配置。}}
  \label{tab:falcon-config}
  \begin{tabular}{l|c}
    \hline
    \textbf{配置项} & \textbf{数值} \\ \hline
    节点数量 & 19 \\
    每节点 CPU & AMD 2$\times$7H12 64-core @ 2.6GHz \\
    每节点内存 & 512 GB \\
    顶点数量 & 1000 万 \\
    自由度 & 3000 万 \\
    目标特征模态数 & 3000 \\
    \hline
  \end{tabular}
\end{table}

在该实验中，3000 阶模态的最大残差为 $1.6\times10^{-4}$，多数高频模态残差接近 $10^{-6}$。此外，与 SLEPc 的对比中，Cluster Pardiso 配置运行超过 7 小时仍未完成；MUMPS 配置在近 2 小时内仅得到 949 阶（阈值 $3\times10^{-2}$）。

\subsection{积分误差分析的实验细节}
\label{app:integrated-error-exp}

本附录给出正文中积分误差数值实验的完整方法细节。该实验用于比较频率适配基与相位适配基在平均意义下的近似性能。

\begin{figure}[H]
  \centering
  \includegraphics[width=\linewidth]{frequency-vs-phase/integration.pdf}
  \caption{\textbf{整体区域与积分子域。} 通过将子域 $[a,b]$ 中的 $a$ 从 $0$ 变化到 $L-l$ 进行积分，消除由特定子域选择带来的偏置。}
  \label{fig:integration}
\end{figure}

\subsubsection{实验目标}
实验目标是在位置无关的意义下，稳健评估局部基的近似效率。具体地，通过在全域内滑动分析窗口（子域）并对 $L_2$ 误差积分，消除单一固定划分可能引入的特殊共振效应，从而揭示两类基函数的真实基础性能差异。

\subsubsection{实验方法}
\paragraph{全局问题}
目标函数为区间 $[0,L]$ 上带 Dirichlet 边界条件的一维拉普拉斯全局特征模态：
$f(x; n_{\text{global}})=\sin(n_{\text{global}}\pi x/L)$。

\paragraph{滑动子域}
在区间 $[a,a+l]$ 上定义长度固定为 $l$ 的子域。令 $a$ 从 $0$ 变化到 $L-l$，从而在全域内滑动窗口，如 \Cref{fig:integration} 所示。

\paragraph{基函数集合}
给定内部模态数 $M$，在区间 $[a,a+l]$ 上构造两组基。每组总函数数均为 $M+2$：
\begin{itemize}
  \item \textbf{频率适配基：}由 $M$ 个递增频率的正弦模态
  $\{\sin(m\pi(x-a)/l)\}_{m=1}^{M}$ 与 2 个线性界面（Lagrange）模态组成。
  \item \textbf{相位适配基：}由混合的正弦/余弦模态
  $\{\sin(m\pi(x-a)/l),\cos(m\pi(x-a)/l)\}_{m=1}^{\lfloor M/2 \rfloor}$ 组成；当 $M$ 为奇数时再加入一个正弦模态，并同样加入 2 个线性界面模态。
\end{itemize}

\paragraph{误差计算}
对每个窗口位置 $a$，将目标函数 $f(x)$ 投影到所选基函数集上，得到 $L_2$ 近似误差 $E(a)$。热图中汇报的最终指标为\textbf{积分平均误差}：
\begin{equation}
\bar{E}(n_{\text{global}}, M) = \frac{1}{L-l} \int_{0}^{L-l} E(a; n_{\text{global}}, M) \, da.
\end{equation}

\subsubsection{参数设置}
用于生成正文积分误差结果的参数为：
\begin{itemize}
  \item 全局区域长度：$L=10.0$；
  \item 滑动子域长度：$l=2.0$；
  \item 全局模态最高阶：$n_{\text{global,max}}=8$；
  \item 内部模态最大数：$M_{\text{max}}=12$。
\end{itemize}
所有积分均采用高精度 \texttt{scipy.integrate.quad} 数值积分实现。

\subsection{可扩展性基准的代价分解}
\label{sec:breakdown}
强扩展与弱扩展实验的分阶段代价分解见下图。

\begin{figure}[H]
  \centering
  \includegraphics[width=0.85\linewidth]{scaling_aggregation_new/strong_scaling_plots/6_strong_scaling_combined_pie_breakdown.png}
  \caption{\textbf{强扩展时间代价分解。} 预处理阶段负责任务分析与分发；后处理阶段负责结果收集与组装。}
  \label{fig:strong-scaling-breakdown}
\end{figure}

\begin{figure}[H]
  \centering
  \includegraphics[width=0.85\linewidth]{scaling_aggregation_new/weak_scaling_plots/6_weak_scaling_combined_pie_breakdown.png}
  \caption{\textbf{弱扩展时间代价分解。} 预处理阶段负责任务分析与分发；后处理阶段负责结果收集与组装。}
  \label{fig:weak-scaling-breakdown}
\end{figure}
